\section{Cyclotomic Fields}

\begin{definition}[\textbf{Primitive  n-th root}]
	Suppose \( n > 0 \) and \( \mu_n \subseteq \CC^\times \), and
	\[
		\begin{aligned}
			\mu_n & = \{\zeta \in \CC^\times \; | \; \zeta^n = 1\}                                  \\
			      & = \{\cos(\frac{2k\pi}{n}) + i\sin(\frac{2k \pi}{n}) \; | \; 0 \leq k \leq n-1\}
		\end{aligned}
	\]
	being a cyclic group. A \underline{primitive \( n \)-th root} of \( 1 \) is \( \zeta \) is a
	generator of this group.
\end{definition}
The goal is to understand the field extension \( \QQ \hookrightarrow \QQ(\mu_n) = \QQ(\zeta) \)
where \( \zeta \) is the primitive \( n \)-th root of the splittiing field of \( x^n-1 \).

\begin{definition}[\textbf{n-th Cyclotomic Polynomial}]
	The \( n \)-th cyclotomic polynomial \( \Phi_n \) is given by
	\[
		\Phi_n = \prod_{\zeta = n\text{-th primitive root of 1}} (x-\zeta) \in \CC[x]
	\]
	which is monic.
\end{definition}

\begin{note}
	If \( \zeta \) is the \( n \)-th primitive root of \( 1 \), then the primitive \( n \)-th root of
	\( 1 \) is given by:
	\[
		\{\zeta^l \; | \; 1 \leq l \leq n, \gcd(l,n) = 1\}
	\]
	which can be computed by Euler's formula. In particular:
	\[
		\deg (\Phi_n) = \vp(n) = \# \{l \; | \; 1\leq l \leq n, \gcd(l,n) = 1\}
	\]
\end{note}

\begin{eg}
	If \( p > 0 \) is prime, then the degree of \( \Phi_n = p - 1 \), and we have:
	\[
		\Phi_n (x) = \frac{x^p - 1}{x - 1} = x^{p-1} + \cdots + x - 1 \in \ZZ[x]
	\]
\end{eg}

\begin{proposition}[\textbf{General Properties of \( \Phi_n \)}]
	\leavevmode
	\begin{enumerate}
		\item \( x^n - 1 = \prod_{d \mid n} \Phi_d(x) \)
		      \begin{proof}
			      For \( x^n - 1 \) there is clearly no multiple roots. See that
			      \[
				      x^n - 1 = \prod_{\zeta \in \mu_n} (x - \zeta)
			      \]
			      with
			      \[
				      \mu_n = \bigsqcup_{d\mid n} \{\text{primitive } d \text{-th root of 1}\}
			      \]
		      \end{proof}
		\item \( \Phi_n \in \ZZ[x] \; \forall \; n \geq 1 \)
		      \begin{proof}
			      Argue by induction on \( n \geq 1 \).
			      \begin{itemize}
				      \item \( n = 1 \): \( \Phi_1 = x - 1 \) which is ok.
				      \item Inductive step:
				            \[
					            x^n  -1 = \prod_{d\mid n} \Phi_d(x) \implies \Phi_n(x) \cdot \underbrace{\prod_{d\mid n, n\ne
							            n} \Phi_d(x)}_{\in \ZZ[x] \text{ by induction}} = x^n -1 \in \ZZ[x]
				            \]
				            Then it is enough to show the following:
				            \begin{proposition}
					            Let \( S \) be commutative ring, and \( R \subseteq S \) being a subring, if
					            \( f,g,h \) are monic polynomials in \( S[x] \), such that \( f=gh \) and \(
					            f,g \in R[x] \), then \( h \in R[x] \).
				            \end{proposition}

				            \begin{proof}
					            Denote:
					            \[
						            \begin{aligned}
							            g & = x^m + b_1 x^{m-1} + \cdots + b_m \\
							            h & = x^n + c_1 x^{n-1} + \cdots + c_n
						            \end{aligned}
					            \]
					            See that the coefficient of \( x^{i + m + n} \) in \( f \) lies in \( R \) for
					            \( 1 \leq i \leq n \), which is actually
					            \[
						            c_i + \underbrace{c_{i-1}b_1 + \cdots + c_1 b_{i-1} + b_i}_{\in R \text{ by
								            induction on $i$, since \( b_j \in R \; \forall \; j \)}} \implies c_i \in R
					            \]
				            \end{proof}
			      \end{itemize}
		      \end{proof}
		\item \( \Phi_n(x) \) is irreducible ($\implies \Phi_n(x)$ is the minimal polynomial of any
		      primitive \( n \)-th root of \( 1 \))
		      \begin{proof}
			      Let \( \zeta \) be the primitive \( n \)-th root of \( 1 \) and let \( f\in \QQ[x] \) be
			      the minimal polynomial of \( \zeta \) over \( \QQ \). Since \( \Phi_n(\zeta) = 0
			      \implies f \mid \Phi_n \). Now let \( \Phi_n = f g \) for some \( g \in \QQ[x] \), then
			      \( f, \Phi_n \) is monic \( \implies g \) is monic.
			      \begin{claim}
				      \( f,g \in \ZZ[x] \).
			      \end{claim}
			      \begin{proof}
				      Recall that if \( P \in \QQ[x] \) is any non-zero polynomial, then \( P = c(P) \cdot P' \) with
				      \( c(P) \in \QQ \) and \( P' \in \ZZ[x] \) being primitive. Thus
				      \[
					      1 = c(\Phi_n) \sim c(f) c(g) \implies 1 = c(\Phi_n) = \pm  c(f) \cdot c(g)
				      \]
				      \begin{note}
					      \( f,g \) being monic \( \implies \frac{1}{c(f)} \cdot \frac{1}{c(g)} \in \ZZ \)
				      \end{note}
				      If \( a \in \ZZ_{>0} \), s.t. \( a f \in \ZZ[x] \), then \( ac(f) = c(af) =
				      \gcd(\text{coefficient of } af) \mid a \implies \exists \; q \in \ZZ \), s.t. \( a c(f) \cdot q = a \implies
				      c(f) = \frac{1}{q}\), hence \( c(f), c(g) = \pm 1 \implies f,g \in \ZZ[x] \).
			      \end{proof}

			      \begin{claim}
				      Fix prime \( p \), s.t. \( p\nmid n \). If \( \alpha \) is a root of \( f \implies
				      \alpha^p \) is a root of \( f \). This implies that \( f = \Phi_n \), and hence \( \Phi_n \) is
				      irreducible. Why is that? Given our root \( \zeta \) of \( f \), every root of \( \Phi_n \) is
				      of the form \( \zeta^m \) where \( \gcd(m,n) = 1 \).

				      Now if \( m = p_1\cdots p_k \) where \( p_i \) is prime, this claim + induction on \( i \) implies that
				      \( \zeta^{p_1\cdots p_i} \) are root of \( f \) for \( 1\leq i \leq k \implies \zeta^m \) is root
				      of \( f \). This holds for all primitive \( n \)-th root of \( 1 \implies f = \Phi_n \), and hence
				      \( \Phi_n \) is irreducible.
			      \end{claim}

			      \begin{proof}
				      We use the ring homomorphism:
				      \[
					      \begin{aligned}
						      \ZZ[x]           & \to \bace{\qo{\ZZ}{p\ZZ}}[x]                   \\
						      x                & \mapsto x                                      \\
						      \sum a_i x^i = Q & \mapsto \overline{Q} = \sum \overline{a_i} x^i
					      \end{aligned}
				      \]
				      We argue by contradiction, suppose \( \alpha \) is a root of \( f \), such that \(
				      \alpha^p  \ne \) root of \( f  \implies  g(\alpha^p) = 0\) by the fact that \( \Phi_n
				      = fg\). Thus
				      \[
					      f(x) \mid g(x^p) \text{ in } \QQ[x] (f \text{ irreducible } \implies f \text{
						      minimal polynomial of } \alpha^p)
				      \]
				      by the fact that \( g(x^p) = f(x) \cdot h(x), h(x) \in \QQ[x] \implies \) by monic \(
				      \implies h \in \ZZ[x] \). Now \( \Phi_n = fg \implies \) In \( \qo{\ZZ}{p\ZZ}[x] \):
				      \( \overline{\Phi_n}(x) = \overline{f}(x) \cdot \overline{g}(x) \), thus
				      \[
					      \begin{aligned}
						      (\overline{\Phi_n}(x))^p = \overline{\Phi_n}(x^p) & = \overline{f} (x^p) \cdot
						      \overline{g}(x^p)                                                                   \\
						                                                        & = \overline{f}(x^p) \cdot
						      \overline{f} (x) \cdot
						      \overline{h} (x)                                                                    \\
						                                                        & = (\overline{f}(x))^p \cdot
						      \overline{f}(x) \cdot
						      \overline{h}(x)                                                                     \\
						                                                        & = (\overline{f}(x))^{p+1} \cdot
						      \overline{h}(x)
					      \end{aligned}
				      \]
				      where we have that \( \overline{f} (x^p) = (\overline{f}(x))^p \) since \( a^p = a \;
				      \forall \; a \in \qo{\ZZ}{p\ZZ}\) and \( u \mapsto u^p \) is a ring homomorphism.
				      \begin{note}
					      \( \overline{\Phi_n} \) doesn't have multiple roots.
				      \end{note}
				      Since \( \overline{\Phi_n} \mid x^n - 1 \) in \( \qo{\ZZ}{p\ZZ} \) and since \( (x^n -
				      1)' = nx^{n-1} \ne 0 \) has only root \( x = 0 \) and \( x = 0 \) not root of \( x^n -
				      1 \implies \overline{\Phi_n}\) has only simple root.

				      Hence, all roots of \( (\overline{\Phi_n})^p \) have multiplicity being \( p \), but
				      \( (\overline{f})^{p+1} \) has roots of multiplicity \( \geq p + 1 \), leading to
				      contradiction \( \lightning \).
			      \end{proof}
		      \end{proof}
	\end{enumerate}
\end{proposition}

\begin{conclusion}
	\leavevmode
	\begin{itemize}
		\item \( [\QQ(\zeta) : \QQ] = \vp(n) \) and the minimal polynomial of \( \zeta \) is \( \Phi_n(x) \), where
		      \( \zeta \) is the primitive \( n \)-th root of \( 1 \).
		\item \( \QQ \hookrightarrow \QQ(\zeta) \) is splitting field of \( x^n - 1 \), and it is
		      normal, separable(by that char\( (\QQ) = 0\)) \( \implies \) it is a finite Galois extension.

		      Then what iis the Galois group?
		      \[
			      \begin{aligned}
				      G \bace{\qo{\QQ(\zeta)}{\QQ}} & \xrightarrow{\sim} \bace{\qo{\ZZ}{n\ZZ}}^\times \\
				      \sigma                        & \mapsto l                                       \\
				      \text{where } \sigma(\zeta)   & = \zeta^l \text{ for some } l, \gcd(l,n) = 1
			      \end{aligned}
		      \]
		      this is a group homomorphism, it is injective since \( \zeta \) generates \( \QQ(\zeta) \)
		      over \( \QQ \). Now both sides have \( \vp(n) \) elements, and therefore being an isomorphism.

		      In particular, the galois group \( G \) is abelian. Thus \( \forall \; \QQ \hookrightarrow
		      K \hookrightarrow \QQ(\zeta)\) with the first extension being normal, see that \(
		      G(\qo{K}{\QQ}) \) is abelian.
	\end{itemize}
\end{conclusion}

\begin{theorem}[\textbf{Kroneckr-Weber Theorem}]
	If \( \qo{K}{\QQ} \) is finite Galois extension s.t. \( G(\qo{K}{\QQ}) \) is abelian, then there
	exists \( n \), s.t. \( K \hookrightarrow \QQ(\mu_n) \).
\end{theorem}

\begin{remark}[\textbf{Open problem}]
	Is it true that every finite group \( G \) is isomorphic to \( G(\qo{K}{\QQ}) \) for some finite
	Galois extension \( \QQ \hookrightarrow K \)?
\end{remark}

\begin{note}
	We can use the description of Galois groups of cyclotomic field to treat the case \( G \) is abelian.

	Suppose
	\[
		G \cong \qo{\ZZ}{n_1\ZZ} \times \cdots \times \qo{\ZZ}{n_r\ZZ}
	\]
	choose \( p_1,\ldots, p_r \) being distinct primes, s.t. \( p_i \equiv 1 \pmod{n_i} \; \forall \; i \), this is a
	consequence of Dirichlet's Theorem about primes in arithmetic progression. Take \( n = p_1\cdots p_r \), by the
	Chinese Remainder Theorem, we have:
	\[
		\qo{\ZZ}{n\ZZ} \cong \prod_{i=1}^r \qo{\ZZ}{p_i \ZZ} \rsa \bace{\qo{\ZZ}{n\ZZ}}^\times \cong
		\prod_{i=1}^r \underbrace{\bace{\qo{\ZZ}{p_i \ZZ}}^\times}_{\text{cyclic, }\cong \qo{\ZZ}{(p_i-1)\ZZ}} \text{ be group isomorphism.}
	\]
	thus there exists a surjective morphism \( (\qo{\ZZ}{n\ZZ})^\times \twoheadrightarrow G \implies
	\exists \QQ \hookrightarrow K \hookrightarrow \QQ(\mu_n) \) where the first extension is Galois,
	such that
	\[
		G\bace{\qo{K}{\QQ}} \cong G
	\]
\end{note}

